% Font "FandolSong-Regular" used by ctex does not have GSUB table.
% It causes warnings that do not affect compilation and can be ignored.
\PassOptionsToPackage{quiet}{fontspec}

\documentclass[openany, 10pt]{book}

%%%%%%%%%%%%%%%%%%%%%%%%%%%%%%%%%%%%%%%%%%%%%%%%%%%
%%% Find  10-readme-for-trm-repo.tex at:        %%%
%%% https://www.overleaf.com/read/bwdjvfjksvjy/ %%%
%%%%%%%%%%%%%%%%%%%%%%%%%%%%%%%%%%%%%%%%%%%%%%%%%%%

% >>> R?: marks a repo project setting number or important notice

% >>> R1: Configure chip-specific settings <<<
\usepackage{chip-spec-settings}

% Imports packages and global formatting settings
\usepackage{preamble-trm-repo}

% Allow Chinese typesetting
\usepackage{ctex}

% >>> R2: In file below, comment out the modules already included into
% this project to avoid warnings for multiple label definitions <<<
\myexternaldocument{temp-labels-cn}

% Make text left-aligned and allow latex to break pages naturally, comment out for full justification
\raggedright
\raggedbottom

% Keep commented out!
% >>> Only uncomment to generate the list of added hypertargets <<<
%\toggletrue{addedhypertargets}


\begin{document}

\selectlanguage{chinese}
\setchineseformatting

% Initial checks make sure that all variables for cross-project shared
% content are configured correctly. If not, guidance will be provided.
%%%%%%%%%%%%%%%%%%%%%%%%%%%%%%%%%%%%%%%%%%%%%%%%%%%%%%%%%%%%%%%%%%%%%%%%%%
%%% For complete instructions:                                         %%%
%%% 1. Go to https://www.overleaf.com/project/60e2ed1aefb70e595778b04b %%%
%%% 2. Compile readme-general.tex                                      %%%
%%% 3. See Section "Initial Checks"                                    %%%
%%%%%%%%%%%%%%%%%%%%%%%%%%%%%%%%%%%%%%%%%%%%%%%%%%%%%%%%%%%%%%%%%%%%%%%%%%


%%% Warning for projects with undefined docprojtype %%%

% Check if the variable docprojtype has been added to the project
% to mark the project type for conditional inclusion of content, etc.
\ifcsname docprojtype\endcsname%

    % Check if the variable docprojtype is assigned a defined document project name
    \ifthenelse{%
      \equal{\docprojtype}{chip-datasheet} \OR
      \equal{\docprojtype}{trm} \OR
      \equal{\docprojtype}{hdg} \OR
      \equal{\docprojtype}{mod-datasheet} \OR
      \equal{\docprojtype}{errata}
    }{%
    % The variable docprojtype has been assigned a defined document project name
    % Proceed compilation.
    }{%
    % The variable docprojtype has NOT been assigned a defined document project name
    \begin{tiplisting}
      In your \texttt{preamble-....sty}, assign one of the following values to \texttt{docprojtype}, depending on your document type:
      \begin{itemize}
        \item \texttt{chip-datasheet}
        \item \texttt{trm}
        \item \texttt{hdg}
        \item \texttt{mod-datasheet}
        \item \texttt{errata}
      \end{itemize}
    \end{tiplisting}
    \errmessage{The variable 'docprojtype' has an unauthorized value}
    \stop
    }%

\else
    % The variable docprojtype has not been added to the project
    \begin{tiplisting}
      The code block below defines the variable \texttt{docprojtype} to mark document project type for conditional inclusion of content, etc.
      \begin{enumerate}
        \item \textbf{Insert the code block below} into your \texttt{preamble-....sty}.
        \item \textbf{Assign one of the following values} to \texttt{docprojtype}, depending on your document type:
	    \begin{itemize}
	        \item \texttt{chip-datasheet}
	        \item \texttt{trm}
	        \item \texttt{hdg}
	        \item \texttt{mod-datasheet}
	        \item \texttt{errata}
	    \end{itemize}
      \end{enumerate}

	    Example:

\begin{lstlisting}[numbers=none]
%%%%%%%%%%%%%%%%%%%%%%%%
%%% System Variables %%%
%%%%%%%%%%%%%%%%%%%%%%%%

% >>> Variable to mark document project type
%      for conditional inclusion of content, etc. <<<
\newcommand\docprojtype{<assign doc type>}
\end{lstlisting}

    \end{tiplisting}

    \errmessage{Define the variable 'docprojtype'}
    \stop

\fi% \docprojtype

\begin{titlepage}


\newgeometry{left=0cm,right=0cm,bottom=0cm}

\begin{flushleft}
{\color{myred}\rule{30pt}{2.7cm}}
    \hspace{0.2cm}
    \begin{minipage}[b]{18cm}
    {\fontsize{36}{48}\textbf{\chipname{}}}
    \vspace{\baselineskip}
    \vfill
    {\fontsize{28}{18}\textbf{\color{mygrayy}\GetTranslation{doctype}}}\smallskip
    \hspace{0.2cm}
    {\fontsize{18}{8}\textbf{\color{mygrayy}\GetTranslation{docversionprefix}\docversion{}}}
    \end{minipage}
\end{flushleft}

\vspace{48em}

\begin{flushright}
    \begin{tikzpicture}[remember picture,overlay]
        \node[anchor=north west, 
        xshift=16.5cm, 
        yshift=5.5cm]
        at (current page.south west) 
        {\includegraphics[height=2.2cm]{esp-logo-standard-vertical.pdf}};
    \end{tikzpicture}
    \begin{tikzpicture}[remember picture,overlay]
        \node[anchor=north west, 
            xshift=20cm, 
            yshift=5.5cm] 
            at (current page.south west) 
            {\color{myred}\rule{30pt}{2.2cm}};
    \end{tikzpicture}    
\end{flushright}

\vspace{9em}

\centerline{www.espressif.com}


\end{titlepage}

\restoregeometry
\clearpage

\clearpage
\thispagestyle{empty}

\setcounter{page}{2}

\subsection*{关于本文档}

%\subsection*{Preliminary Release Note}


{\bfseries \chipname{} \GetTranslation{doctype}}面向使用 \chipname{} 系列产品进行底层软件开发的人员，介绍了 \chipname{} 系列产品中内置的硬件模块，包括概述、功能列表、硬件架构、编程指南、寄存器列表等信息。

\subsection*{本文档中的跳转}
在本文档中实现跳转，请参考以下建议：
\begin{itemize}
    \item \hyperref[tab:release-status]{发布进度速览}（下一页）罗列了本文档中的所有章节，您可以从这里快速跳转至某个具体章节。
    \item 您还可以通过文档左侧的\textbf{书签}，从文中的任何位置直接跳转至另一个章节。注意，本文档已设置默认打开\textbf{书签}功能，但一些 PDF 阅读器或浏览器会忽略此设置。因此，如果您无法找到\textbf{书签}功能，请尝试以下方法：
    \begin{itemize}
        \item 在您的浏览器中安装 PDF 阅读器拓展；
        \item 下载本文档，使用本地 PDF 阅读器进行浏览；
        \item 配置您的 PDF 阅读器，使其默认打开\textbf{书签}功能。
    \end{itemize}
    \item 大多数 PDF 阅读器均支持跳转功能，允许您借助按钮、菜单选项或快捷键进行跳转（\textbf{向上}、\textbf{向下}、\textbf{向前}
    、\textbf{向后}、\textbf{后退}、\textbf{前进}
    及\textbf{前往页面}）等。
    \item 此外，您还可以使用本文档内置的 \textbf{GoBack} 按钮（每页右上角）快速后退至跳转之前的位置。注意，本功能仅适用于 Acrobat 系列的 PDF 阅读器（比如 Acrobat Reader 和 Adobe DC）以及内置 Acrobat 系列 PDF 阅读器或拓展的浏览器（比如 Firefox）。
\end{itemize}

\clearpage
% % \thispagestyle{empty}

\setcounter{page}{3}

\subsection*{发布进度速览}
\markboth{\it{发布进度速览}}{\Navigation}

\begin{table}[H]
\centering
\label{tab:release-status}
\begin{tabular}{|c|L{5cm}|c|c|L{5cm}|c|}
\hline
\rowcolor{lightgray}
\textbf{No.} & \textbf{章节} & \textbf{最新进度}  & \textbf{No.} & \textbf{章节} & \textbf{最新进度} \\ \hline
\multicolumn{3}{|c|}{\textbf{第 \uppercase\expandafter{\romannumeral1} 卷：微处理器和主机}} & 20    &  \nameref*{mod:rsa}    &    \hyperref[mod:rsa]{已发布} \\\hline
1    &  \nameref*{mod:pie}    & \hyperref[mod:pie]{已发布}   & 21    &   \nameref*{mod:hmac}    &    \hyperref[mod:hmac]{已发布}\\\hline
2    &  \nameref*{mod:ulp}    &    \hyperref[mod:ulp]{已发布}   & 22    &  \nameref*{mod:digsig}    &    \hyperref[mod:digsig]{已发布}     \\\hline
3    &  \nameref*{mod:dma}    &    \hyperref[mod:dma]{已发布}   & 23    &  \nameref*{mod:extmemencr}    &    \hyperref[mod:extmemencr]{已发布}  \\\hline
\multicolumn{3}{|c|}{\textbf{第 \uppercase\expandafter{\romannumeral2} 卷：存储器组织结构}} & 24    &  \nameref*{mod:glitch}    &    \hyperref[mod:glitch]{已发布} \\\hline
4    &  \nameref*{mod:sysmem}    &    \hyperref[mod:sysmem]{已发布}  & 25    &    \nameref*{mod:rng}    &    \hyperref[mod:rng]{已发布} \\\hline
5    &  \nameref*{mod:efuse}    &    \hyperref[mod:efuse]{已发布}  & \multicolumn{3}{|c|}{\textbf{第 \uppercase\expandafter{\romannumeral5} 卷：通讯接口}}    \\\hline
\multicolumn{3}{|c|}{\textbf{第 \uppercase\expandafter{\romannumeral3} 卷：系统组件}} & 26    &  \nameref*{mod:uart}    &    \hyperref[mod:uart]{已发布} \\\hline
6    &  \nameref*{mod:iomuxgpio}    &    \hyperref[mod:iomuxgpio]{已发布}   & 27    &  \nameref*{mod:i2c}    &    \hyperref[mod:i2c]{已发布}  \\\hline
7    &  \nameref*{mod:resclk}    &    \hyperref[mod:resclk]{已发布}  & 28    &  \nameref*{mod:i2s}    &    \hyperref[mod:i2s]{已发布}    \\\hline
8    &  \nameref*{mod:bootctrl}    &    \hyperref[mod:bootctrl]{已发布} & 29    &   \nameref*{mod:camlcdctrl}    &    \hyperref[mod:camlcdctrl]{已发布}  \\\hline
9    &  \nameref*{mod:intmtrx}    &    \hyperref[mod:intmtrx]{已发布} & 30    &  \nameref*{mod:spi}    &    \hyperref[mod:spi]{已发布}   \\\hline
10    &  \nameref*{mod:lowpowm}    &    \hyperref[mod:lowpowm]{已发布}  & 31    &  \nameref*{mod:twai}    &    \hyperref[mod:twai]{已发布}  \\\hline
11    &  \nameref*{mod:systimer}    &    \hyperref[mod:systimer]{已发布}  & 32    &    \nameref*{mod:usbotg}    &    \hyperref[mod:usbotg]{已发布} \\\hline
12    &  \nameref*{mod:timg}    &    \hyperref[mod:timg]{已发布}  & 33    &  \nameref*{mod:usbserialjtag}    &    \hyperref[mod:usbserialjtag]{已发布}    \\\hline
13    &  \nameref*{mod:wdt}    &    \hyperref[mod:wdt]{已发布}   & 34    &  \nameref*{mod:sdhost}    &    \hyperref[mod:sdhost]{已发布}  \\\hline
14    &  \nameref*{mod:xtwdt}    &    \hyperref[mod:xtwdt]{已发布}   & 35    &  \nameref*{mod:ledpwm}    &    \hyperref[mod:ledpwm]{已发布}   \\\hline
15    &  \nameref*{mod:permctrl}    &    \hyperref[mod:permctrl]{已发布}  & 36    &  \nameref*{mod:mcpwm}    &    \hyperref[mod:mcpwm]{已发布}  \\\hline
16    &  \nameref*{mod:wctrl}    &    \hyperref[mod:wctrl]{已发布}  & 37    &  \nameref*{mod:rmt}    &    \hyperref[mod:rmt]{已发布}   \\\hline
17    &  \nameref*{mod:sysreg}    &    \hyperref[mod:sysreg]{已发布}  & 38   & \nameref*{mod:pcnt}    &    \hyperref[mod:pcnt]{已发布}  \\\hline
\multicolumn{3}{|c|}{\textbf{第 \uppercase\expandafter{\romannumeral4} 卷：加密和安全组件}} & 39    &  \nameref*{mod:sensor}    &    \hyperref[mod:sensor]{已发布}   \\\hline
18    &  \nameref*{mod:sha}    &    \hyperref[mod:sha]{已发布}     & 38   & \nameref*{mod:pcnt}    &    \hyperref[mod:pcnt]{已发布}  \\\hline
19    &  \nameref*{mod:aes}    &    \hyperref[mod:aes]{已发布}  & \multicolumn{3}{|c|}{\textbf{第 \uppercase\expandafter{\romannumeral6} 卷：模拟信号处理}} \\\hline
20    &  \nameref*{mod:rsa}    &    \hyperref[mod:rsa]{已发布}  & 39    &  \nameref*{mod:sensor}    &    \hyperref[mod:sensor]{已发布}   \\\hline
\end{tabular}
\end{table}


\urlstyle{rm}

\begin{tiplisting}
    \vspace{-3.7em}
    \begin{longtable}[l]{ m{13.8cm}  p{2.5cm} }

        % table contents
        \makecell[l]{点击链接或扫描二维码确保您使用的是最新版本的文档：\\
            %\vspace{0.5em}\\
            %\href{\docpathlatestCN}{\docpathlatestCN}}
            \parbox{13.8cm}{\url{\docpathlatestCN}}}
        & \qrcode[height=2cm]{\docpathlatestCN}\\
    \end{longtable}
\end{tiplisting}

\setcounter{table}{0}



\tableofcontents
\clearpage
\listoftables
\clearpage
\listoffigures

% >>> R3: List all todo notes from the current project? <<<
% 'Yes' - uncomment the lines below; 'No' - comment out
% Don't forget to deactivate before publishing
%\todototoc
%\listoftodos
%{\let\clearpage\relax \listofdonetodos}


% Included modules
% Uncomment to include


%%%%%%%%%%%%%%%%%%%%%%%%%%%%
%%% Module Group: System %%%
%%%%%%%%%%%%%%%%%%%%%%%%%%%%

% \part{微处理器和主机}

% % Note: Updates may be necessary based on the specific SoC. For instance, Xtensa-based SoCs lack CPU descriptions, and certain SoCs may omit processor instruction extensions.

% 该部分介绍构成芯片系统的基本要素，包括超低功耗协处理器架构、CPU 的扩展指令、直接内存访问 (DMA) 控制器等。

% \subfileinclude{08-PIE__CN}
% \subfileinclude{26-ULP__CN}
% \subfileinclude{16-DMA__CN}

% \part{存储器组织结构}

% 该部分详细描述系统的存储器组织结构，例如 RAM、ROM、eFuse 和外部存储器的组织和映射，为理解与存储器相关的子系统提供了框架。

% \subfileinclude{09-SYSMEM__CN}
% \subfileinclude{07-eFuse__CN}

% \part{系统组件}

% % Note: Tailor the description based on the specific components present in the SoC.

% 该部分介绍芯片的系统级功能，包括系统启动、时钟、GPIO、定时器、看门狗、中断处理、低功耗管理以及各种系统寄存器。

% \subfileinclude{21-IOMUXGPIO__CN}
% \subfileinclude{02-RESCLK__CN}
% \subfileinclude{35-BOOTCTRL__CN}
% \subfileinclude{22-INTMTRX__CN}
% \subfileinclude{28-LOWPOWM__CN}
% \subfileinclude{34-SYSTIMER__CN}
% \subfileinclude{14-TIMG__CN}
% \subfileinclude{15-WDT__CN}
% \subfileinclude{37-XTWDT__CN}
% \subfileinclude{10-PERMCTRL__CN}
% \subfileinclude{38-WCTRL__CN}
% \subfileinclude{23-SYSREG__CN}

% \part{加密和安全组件}

% 该部分专门介绍芯片的安全功能，包括 SHA、AES 等加密加速器，数字签名、随机数生成、加密/解密算法等安全特性，展示了芯片在加密和安全数据处理方面的性能。

% \subfileinclude{27-SHA__CN}
% \subfileinclude{30-AES__CN}
% \subfileinclude{31-RSA__CN}
% \subfileinclude{29-HMAC__CN}
% \subfileinclude{13-DIGSIG__CN}
% \subfileinclude{33-EXTMEMENCR__CN}
% \subfileinclude{36-GLITCH__CN}
% \subfileinclude{06-RNG__CN}


%%%%%%%%%%%%%%%%%%%%%%%%%%%%%%%%%
%%% Module Group: Peripherals %%%
%%%%%%%%%%%%%%%%%%%%%%%%%%%%%%%%%


\part{通讯接口}

% Note: Update the description based on the actual components implemented in the SoC, e.g., some SoCs do not have RMT.

该部分涉及系统的连接性，介绍各种带有通讯接口的模块，如 I2C、I2S、SPI、UART、USB 等，同时还包括其他用于遥控、电机控制、LED 控制等功能的信号生成接口。

% \subfileinclude{12-UART__CN}
\subfileinclude{01-I2C__CN}
\subfileinclude{03-I2S__CN}
% \subfileinclude{04-CAMLCDCTRL__CN}
% \subfileinclude{24-SPI__CN}
% \subfileinclude{32-TWAI__CN}
% \subfileinclude{11-USBOTG__CN}
% \subfileinclude{39-USBSERIALJTAG__CN}
% \subfileinclude{18-SDHOST__CN}
% \subfileinclude{19-LEDPWM__CN}
% \subfileinclude{17-MCPWM__CN}
% \subfileinclude{05-RMT__CN}
% \subfileinclude{20-PCNT__CN}

% \part{模拟信号处理}

% % Note: Tailor the description based on the actual components implemented in the SoC.

% 该部分描述了与模数转换、片上传感器以及温度感应等功能相关的模块，展示了系统在处理模拟信号方面的性能。

% \subfileinclude{25-SENSOR__CN}


%%%%%%%%%%%%%%%%
%%% Appendix %%%
%%%%%%%%%%%%%%%%

\part{附录}\label{appendix}

附录部分包括以下内容：

\begin{itemize}
    \item \hyperref[sec:resources]{\textit{相关文档和资源}}
    \item \hyperref[glossary]{\textit{词汇列表}}
    \item \hyperref[programming-reserved-field]{\textit{如何配置寄存器的保留域}}
    \item \hyperref[interrupt-config-registers]{\textit{中断配置寄存器}}
    \item \hyperref[revision-history]{\textit{修订历史}}
\end{itemize}


%%%%%%%%%%%%%%%%%%%%%%%%%%%%%%%%%%%%%%%%%%%
%%% Related Documentation and Resources %%%
%%%%%%%%%%%%%%%%%%%%%%%%%%%%%%%%%%%%%%%%%%%

\clearpage
%%%%%%%%%%%%%%%%%%%%%%%%%%%%%%%%%%%%%%%%%%%%%%%%%%%%%%%%%%%%%%%%%%%%%%%%%%
%%% For complete instructions:                                         %%%
%%% 1. Go to https://www.overleaf.com/project/60e2ed1aefb70e595778b04b %%%
%%% 2. Compile readme-general.tex                                      %%%
%%% 3. See Section "Add Section Related Documentation and Resources"   %%%
%%%%%%%%%%%%%%%%%%%%%%%%%%%%%%%%%%%%%%%%%%%%%%%%%%%%%%%%%%%%%%%%%%%%%%%%%%


% >>> Add chip name here <<<
% Always add '%' after the last chip name below (before the line '}')
\newcommand\definedchipnames{%
ESP32,
ESP32-S2,
ESP32-S3,
ESP32-C3,
ESP32-H2,
ESP8685,
ESP8684,
ESP32-C6,
ESP32-PICO,
ESP32-S3-PICO-1,
ESP8686,
ESP32-P4,
ESP32-C5,
ESP32-C61/ESP32-C31,
ESP-Shelly-C38F,
XMCC3ESP-FH%
}


% Code needed for the Introduced variable \definedchipnames
% to define all chip names in one place
% See https://tex.stackexchange.com/questions/642211/using-a-macro-to-pass-parameters-to-a-command/
% replace \docsvlist in \iftagged with \xdocsvlist
\patchcmd{\iftagged}{\docsvlist}{\xdocsvlist}{}{}
\patchcmd{\tagged}{\docsvlist}{\xdocsvlist}{}{}
\patchcmd{\untagged}{\docsvlist}{\xdocsvlist}{}{}
% \xdocsvlist fully expands its argument before passing it to \docsvlist
\newcommand{\xdocsvlist}[1]{%
  \expandafter\docsvlist\expandafter{\expanded{#1}}%
}

\phantomsection
\chapter*{相关文档和资源}\label{sec:resources}
\markboth{\it{相关文档和资源}}{\Navigation}
\addcontentsline{toc}{chapter}{相关文档和资源}


%%%%%%%%%%%%%%%%%%%%%%%%%%%%%%%%%%%%%%%%%%%%%%%%%
%%% Warning for projects with undefined chips %%%
%%%%%%%%%%%%%%%%%%%%%%%%%%%%%%%%%%%%%%%%%%%%%%%%%

% Show content below if undefined chip name encountered
\untagged{\definedchipnames}
{%
    \begin{tiplisting}
        Assign a defined chip name in your preamble > variable "chipseries"
    
	    Otherwise, define chip in 
	    \begin{itemize}
	        \item \texttt{resources-link-defs}
	        \item \texttt{resources-[doc-type]\_\_EN}
	        \item \texttt{resources-[doc-type]\_\_CN}
	    \end{itemize}
    \end{tiplisting}
}% \untagged{\definedchipnames}



%%%%%%%%%%%%%%%%%%%%%%%%%%%%%%%%%%%%%%%%%%%%%%%
%%% Content for projects with defined chips %%%
%%%%%%%%%%%%%%%%%%%%%%%%%%%%%%%%%%%%%%%%%%%%%%%

% Show content below if defined chip name encountered
\tagged{\definedchipnames}
{%

%%% Related Documentation %%%

\vspace{-0.5em}
\begin{minipage}{\textwidth}

\section*{相关文档}

\begin{itemize}
  % Include this link only if it is defined in resources-link-defs.tex
  \ifcsname linkchipdatasheetCN\endcsname%
    % Exclude this link if this section is added to a chip datasheet project
    \ifthenelse{\equal{\docprojtype}{chip-datasheet}}{}%
    {%
      \item \href{\linkprefix\linkchipdatasheetCN}{《\chipseries{} 技术规格书》} --
        {\small 提供 \chipseries{} 芯片的硬件技术规格。}%
    }%
  \fi%

  % Include this link only if it is defined in resources-link-defs.tex
  \ifcsname linkchiptrmCN\endcsname%
    % Exclude this link if this section is added to a TRM project
    \ifthenelse{\equal{\docprojtype}{trm}}{}%
    {%
      \item \href{\linkprefix\linkchiptrmCN}{《\chipseries{} 技术参考手册》} --
        {\small 提供 \chipseries{} 芯片的存储器和外设的详细使用说明。}%
    }%
  \fi%

  % Include this link only if it is defined in resources-link-defs.tex
  \ifcsname linkchiphwdesignCN\endcsname%
    % Exclude this link if this section is added to a hardware design guidelines project
    \ifthenelse{\equal{\docprojtype}{hdg}}{}%
    {%
      \item \href{\linkprefix\linkchiphwdesignCN}{《\chipseries{} 硬件设计指南》} --
        {\small 提供基于 \chipseries{} 芯片的产品设计规范。}%
    }%
  \fi%

  % Include this link only if it is defined in resources-link-defs.tex
  \ifcsname linkchipecoCN\endcsname%
    \item \href{\linkprefix\linkchipecoCN}{《\chipseries{} 勘误表及解决办法》} --
        {\small 提供关于 \chipseries{} 芯片的设计问题的说明和解决方案。}
  \fi%
  
% Include this link only if it is defined in resources-link-defs.tex
  \ifcsname linkchiperrataCN\endcsname%
   % Exclude this link if this section is added to an errata project
    \ifthenelse{\equal{\docprojtype}{errata}}{}%
    {%
    \item \href{\linkprefix\linkchiperrataCN}{《\chipseries{} 系列芯片勘误表》} --
        {\small 描述 \chipseries{} 系列芯片的已知错误。}
    }%
  \fi%

  % If the variable \productcertificates is defined in a target project
  % append it to the link to filter available certificates
  \ifcsname productcertificates\endcsname%
    \item 证书\\
        \href{\linkprefix\linkcertificatesCN\productcertificates}{\linkprefix\linkcertificatesCN}
  \else
    \PackageWarning{Related Documentation and Resources}{Link extension for certificates not defined}
    \item 证书\\
        \href{\linkprefix\linkcertificatesCN}{\linkprefix\linkcertificatesCN}
  \fi%

  % Include this link only if it is defined in resources-link-defs.tex
  \ifcsname linkpcnCN\endcsname%
    \item \chipseries{} 产品/工艺变更通知 (PCN)\\
        \href{\linkprefix\linkpcnCN}{\linkprefix\linkpcnCN}
  \fi%

  % Include this link only if it is defined in resources-link-defs.tex
  \ifcsname linkadvisoryCN\endcsname%
    \item \chipseries{} 公告 --
        {\small 提供有关安全、bug、兼容性、器件可靠性的信息}\\
        \href{\linkprefix\linkadvisoryCN}{\linkprefix\linkadvisoryCN}
  \fi%

    \item 文档更新和订阅通知\\
        \href{\linkprefix\linkrecentlyupdatedCN}{\linkprefix\linkrecentlyupdatedCN}

\end{itemize}

\end{minipage}


%%% Developer Zone %%%

\begin{minipage}{\textwidth}

\section*{开发者社区}

\begin{itemize}
  % Include this link only if it is defined in resources-link-defs.tex
  \ifcsname linkidfprogguideCN\endcsname%
    \item \href{\linkprefix\linkidfprogguideCN}{《\chipseries{} ESP-IDF 编程指南》} --
        {\small ESP-IDF 开发框架的文档中心。}
  \fi%

    \item ESP-IDF 及 GitHub 上的其它开发框架\\
        \href{\linkprefix\linkgithub}{\linkprefix\linkgithub}

    \item ESP32 论坛 --
        {\small 工程师对工程师 (E2E) 的社区，您可以在这里提出问题、解决问题、分享知识、探索观点。}\\
        \href{\linkprefix\linkespforum}{\linkprefix\linkespforum}
        
    \item \textit{The ESP Journal} --
        {\small 分享乐鑫工程师的最佳实践、技术文章和工作随笔。}\\
        \href{\linkprefix\linkespjournal}{\linkprefix\linkespjournal}

    \item SDK 和演示、App、工具、AT 等下载资源\\
        \href{\linkprefix\linksdkdemoCN}{\linkprefix\linksdkdemoCN}
\end{itemize}

\end{minipage}


%%% Products %%%

\begin{minipage}{\textwidth}

\section*{产品}

\begin{itemize}
    \item \chipseries{} 系列芯片 --
        {\small \chipseries{} 全系列芯片。}\\
        \href{\linkprefix\linkespsocsCN?id=\chipseries}{\linkprefix\linkespsocsCN?id=\chipseries}
    \item \chipseries{} 系列模组 --
        {\small \chipseries{} 全系列模组。}\\
        \href{\linkprefix\linkespmodulesCN?id=\chipseries}{\linkprefix\linkespmodulesCN?id=\chipseries}
    \item \chipseries{} 系列开发板 --
        {\small \chipseries{} 全系列开发板。}\\
        \href{\linkprefix\linkespdevkitsCN?id=\chipseries}{\linkprefix\linkespdevkitsCN?id=\chipseries}
    \item ESP Product Selector（乐鑫产品选型工具）--
        {\small 通过筛选性能参数、进行产品对比快速定位您所需要的产品。}\\
        \href{\linkprefix\linkespproductselectorCN}{\linkprefix\linkespproductselectorCN}
\end{itemize}

\end{minipage}


%%% Contact Us %%%

\begin{minipage}{\textwidth}

\section*{联系我们}

\begin{itemize}
    \item 商务问题、技术支持、电路原理图 \& PCB 设计审阅、购买样品（线上商店）、成为供应商、意见与建议\\
        \href{\linkprefix\linksalesCN}{\linkprefix\linksalesCN}
\end{itemize}

\end{minipage}

}% \tagged{\definedchipnames}


\clearpage
\phantomsection
\chapter*{词汇列表}\label{glossary}
\markboth{\it{词汇列表}}{\Navigation}
\addcontentsline{toc}{chapter}{词汇列表}

\phantomsection
\section*{外设相关词汇}
\addcontentsline{toc}{section}{外设相关词汇}

\begin{longtable}[c]{p{3cm}p{12cm} }

AES & AES 加速器 \\ % 7
BOOTCTRL & 芯片 Boot 控制 \\ %新增
DS & 数字签名 \\ %新增
DMA & DMA 控制器 \\
eFuse & eFuse 控制器 \\
% EMAC & 以太网 MAC
HMAC & HMAC 加速器 \\ %新增
I2C & I2C 控制器 \\
I2S & I2S 控制器 \\
LEDC & LED 控制 PWM \\
MCPWM & 电机控制 PWM \\
% MMU & 存储器管理单元
% MPU & 存储器保护单元
PCNT & 脉冲计数器控制器 \\
% PERI & 外设
% PID & 进程号
RMT & 红外遥控 \\
RNG & 随机数生成器 \\
RSA & RSA 加速器 \\
% RTC & 实时控制器。SOC 中的一组电路，在任何芯片模式下都能随时保持工作
SDHOST & SD/MMC 主机控制器 \\
SHA & SHA 加速器 \\
SPI & SPI 控制器 \\
SYSTIMER & 系统定时器  \\ %新增
TIMG & 定时器组 \\
TWAI & 双线汽车接口 \\ %新增
UART & UART 控制器 \\
ULP 协处理器 & 超低功耗协处理器\\
USB OTG & USB On-The-Go \\ %新增
WDT & 看门狗定时器 \\ 
%XTWDT & XTAL32 看门狗 \\ %新增

\end{longtable}

%%% Registers %%%

% Include sections
% - Abbreviations Related to Registers
% - Access Types for Registers

%%% Abbreviations Related to Registers %%%

\phantomsection
\section*{寄存器相关缩写}\label{glossary-abbr-reg}
\addcontentsline{toc}{section}{寄存器相关缩写}

\begin{longtable}[c]{ R{3cm} p{12cm} }

REG     & \textbf{寄存器}。 \\
SYSREG  & \textbf{系统寄存器}是一组控制系统复位、存储器、时钟、软件中断、电源管理、时钟门控等的寄存器。 \\
ISO     & \textbf{隔离}。如果外设或其他芯片组件断电，其输出信号的管脚（若有）将会浮空。ISO 寄存器会隔离上述引脚并令其保持在某个确定值，以使连接到这些引脚的其他非断电外设/设备免受影响。 \\
NMI     & \textbf{非屏蔽中断}是一种 CPU 指令无法禁用或忽略的硬件中断。出现此类中断说明发生严重错误。 \\
W1TS    & 添加到寄存器/字段名称中的缩写，表示此类寄存器/字段用于置位名称相似寄存器中的相应字段。例如，寄存器 \texttt{GPIO\_ENABLE\_W1TS\_REG} 用于置位寄存器 \texttt{GPIO\_ENABLE\_REG} 中的相应字段。 \\
W1TC    & 与 \textit{W1TS} 相同，但用于清除相应寄存器中的字段。 \\
% FPU   &
% FPD   &
% PLLA  &
% BBPLL &

\end{longtable}


%%% Access Types for Registers %%%

% \clearpage
\phantomsection
\section*{寄存器的访问类型}\label{glossary-access-types}
\addcontentsline{toc}{section}{寄存器的访问类型}

TRM 章节\ 寄存器列表\ 和\ 寄存器\ 详述了寄存器及其字段的访问类型。

常用访问类型及组合如下：

\begin{multicols}{4}
  \begin{itemize}
    \item RO
    \item WO
    \item WT
    \item R/W
    \item WL
    \item R/W/SC
    \item R/W/SS
    \item R/W/SS/SC
    \item R/WC/SS
    \item R/WC/SC
    \item R/WC/SS/SC
    \item R/WS/SC
    \item R/WS/SS
    \item R/WS/SS/SC
    \item R/SS/WTC
    \item R/SC/WTC
    \item R/SS/SC/WTC
    \item RF/WF
    \item R/SS/RC
    \item varies
  \end{itemize}
\end{multicols}

下文提供了所有访问类型的具体描述。


\begin{longtable}[c]{ R{1cm} p{13cm} }

% table contents
R & \textbf{软件可读。} 
	用户软件可以读取此寄存器/字段；通常与其他访问类型结合使用。 \\
RO & \textbf{软件只读。} 
	用户软件只可读取此寄存器/字段。 \\
HRO & \textbf{硬件只读。} 
	仅硬件可以读取此寄存器/字段；用于存储变量参数的默认设置。 \\
W & \textbf{软件可写。} 
	用户软件可以写入此寄存器/字段；通常与其他访问类型结合使用。 \\
WO & \textbf{软件只写。} 
	用户软件只可写入此寄存器/字段。 \\

SS & \textbf{硬件置位。} 
	在指定事件中，硬件自动将 1 写入此寄存器/字段；与一位字段一同使用。 \\
SC & \textbf{硬件清零。} 
	在指定事件中，硬件自动将 0 写入此寄存器/字段；与一位和多位字段一同使用。 \\
SM & \textbf{硬件修改。} 
	在指定事件中，硬件自动将指定值写入此寄存器/字段；与多位字段一同使用。 \\
SU & \textbf{硬件更新。} 
	在指定事件中，硬件自动更新此寄存器/字段；与多位字段一同使用。 \\

RS & \textbf{软件读置位。} 
	如果用户软件读取此寄存器/字段，硬件会自动写 1。 \\
RC & \textbf{软件读清零。} 
	如果用户软件读取此寄存器/字段，硬件会自动写 0。 \\
RF & \textbf{软件读 FIFO。} 
	如果用户软件将新数据写入 FIFO，寄存器/字段会自动读取。 \\

WF & \textbf{软件写 FIFO。} 
	如果用户软件将新数据写入此寄存器/字段，寄存器/字段会自动通过 APB 总线将数据传递到 FIFO。 \\

WS & \textbf{软件写置位。} 
	如果用户软件写入此寄存器/字段，硬件会自动置位此寄存器/字段。 \\
W1S & \textbf{软件写 1 置位。} 
	如果用户软件将 1 写入此寄存器/字段，硬件会自动置位此寄存器/字段。 \\
W0S & \textbf{软件写 0 置位。} 
	如果用户软件将 0 写入此寄存器/字段，硬件会自动置位此寄存器/字段。 \\

WC & \textbf{软件写清零。} 
	如果用户软件写入此寄存器/字段，硬件会自动清零此寄存器/字段。 \\
W1C & \textbf{软件写 1 清零。} 
	如果用户软件将 1 写入此寄存器/字段，硬件会自动清零此寄存器/字段。 \\
W0C & \textbf{软件写 0 清零。} 
	如果用户软件将 0 写入此寄存器/字段，硬件会自动清零此寄存器/字段。 \\

WT & \textbf{软件写产生边沿触发信号。} 
	如果用户软件将 1 写入此字段，将会产生边沿触发信号（APB 总线中的脉冲）或清除相应的 WTC 字段（详见 WTC）。 \\
WTC & \textbf{软件写其他寄存器位清零本寄存器位。} 
	如果用户软件将 1 写入相应的 WT 字段，硬件会自动清除此字段（详见 WT）。 \\

W1T & \textbf{软件写 1 取反。} 
	如果用户软件将 1 写入此字段，硬件会自动取反相应字段，否则不会取反。 \\
W0T & \textbf{软件写 0 取反。} 
	如果用户软件将 0 写入此字段，硬件会自动取反相应字段，否则不会取反。 \\

WL & \textbf{软件仅在锁禁用时写。} 
	如果锁被禁用，用户软件可以写入此寄存器/字段。 \\

varies & \textbf{访问类型不定。} 
	此寄存器中的不同字段访问类型可能不同。 \\

\end{longtable}


% \subsection*{寄存器相关词汇}
% \phantomsection
% \addcontentsline{toc}{subsection}{寄存器相关词汇}

% \begin{longtable}[c]{p{3cm}p{12cm} }
% ISO & 隔离。当模块断电时，其输出的引脚将处于未知状态（某些中间电压）。“ISO” 寄存器将使其输出引脚隔离在一个确定的电压，从而不会影响其他未掉电的工作模块的状态。\\
% NMI & 不可屏蔽中断。 \\
% REG & 寄存器。\\
% R/W & 读/写，软件可读写这些位。 \\
% RO & 只读，软件只能读这些位。 \\
% SYSREG & 系统寄存器。 \\ %新增
% WO & 只写，软件只能写这些位。\\


% \end{longtable}


\clearpage
\phantomsection
\chapter*{如何配置寄存器的保留域}\label{programming-reserved-field}
\markboth{\it{如何配置寄存器的保留域}}{\Navigation}
\addcontentsline{toc}{chapter}{如何配置寄存器的保留域}

\phantomsection
\section*{概述}
\addcontentsline{toc}{section}{概述}

寄存器的保留域指的是寄存器中不对用户开放、或者配置为非默认值时会导致不可预测的结果的域。

\phantomsection
\section*{如何配置保留域}
\addcontentsline{toc}{section}{如何配置保留域}

保留域的值不能修改。由于写寄存器时必须整体写，不能只写部分域，因此，在写带有保留域的寄存器时需要特别注意，只能采用以下两种方式：

\begin{enumerate}
\item 读取寄存器的值，仅修改需要配置的域，然后将修改后的值以及其他未修改的值一起写回寄存器，这样保留域的值就会保持不变。

\noindent 或者

\item 仅修改需要配置的域，然后将保留域的默认值写回寄存器。默认值即寄存器图表中的“Reset”值。例如，\hyperref[fig:registerX]{寄存器 X} 中 Field\_A 的默认值为 1。
\end{enumerate}

\begin{register}{H}{寄存器 X}{地址}\label{fig:registerX}
\regfield{(reserved)}{12}{20}{000000000000}%
\regfield{Field\_C}{4}{16}{{0000}}%
\regfield{(reserved)}{14}{2}{00000000000000}%
\regfield{Field\_B}{1}{1}{{0}}%
\regfield{Field\_A}{1}{0}{{1}}%
\reglabel{Reset}\regnewline%

\end{register}

假设您要将 \hyperref[fig:registerX]{寄存器 X} 的 Field\_A、Field\_B 和 Field\_C 设置为 0x0、0x1 和 0x2，您可以：

\begin{itemize}
\item 使用第 1 种方式，修改这三个域的值，然后把读到的值写回寄存器。假设寄存器读取的值为 0x0000\_0003，修改三个域的值后，将值 0x0002\_0002 写入寄存器。
\item 使用第 2 种方式，修改这三个域的值，然后把保留域的默认值写回寄存器，即把 0x0002\_0002 写入寄存器。
\end{itemize}


\clearpage
\phantomsection
\chapter*{中断配置寄存器}\label{interrupt-config-registers}
\markboth{\it{中断配置寄存器}}{\Navigation}
\addcontentsline{toc}{chapter}{中断配置寄存器}

大部分外设的内部中断源都有以下配置寄存器：

\begin{itemize}
    \item \textbf{RAW}（原始状态）寄存器：该寄存器指示原始中断状态，每个位对应一个内部中断源。当中断源触发时，其 RAW 位为 1。
    
    \item \textbf{ENA}（使能）寄存器：该寄存器用于启用或禁用内部中断源，每个位对应一个内部中断源。
    
    通过操作 ENA 寄存器，可以根据需要屏蔽或取消屏蔽某个内部中断源。当中断源被屏蔽（禁用）时，它不会生成中断信号，但仍可以从 RAW 寄存器中读取其值。
    
    \item \textbf{ST}（状态）寄存器：该寄存器指示中断源的屏蔽状态，每个位对应一个内部中断源。ST 位为 1 代表 RAW 位和 ENA 位都为 1，即中断源已生成且未被屏蔽。RAW 位和 ENA 位的值为其他组合时，ST 位为 0。
    
    ENA/RAW/ST 寄存器的配置见表 \ref{tab:int-raw-st-ena}。

    \item \textbf{CLR}（清除）寄存器：CLR 寄存器负责清除内部中断源。写 1 将清除该位对应的中断源。
\end{itemize}

\begin{longtable}{ | l | l | l | }
\caption{ENA/RAW/ST 寄存器的配置} \label{tab:int-raw-st-ena}
\\ \hline
\rowcolor{lightgray}
\textbf{ENA 位的值} & \textbf{RAW 位的值}  &  \textbf{ST 位的值}\\ \hline
0 & 忽略 & 0 \\ \hline
\multirow{2}*{1} &  0 & 0 \\ \cline{2-3}
 &  1 & 1 \\ \hline
\end{longtable}

\clearpage
\phantomsection
\chapter*{修订历史}\label{revision-history}
\markboth{\it{修订历史}}{\Navigation}
\addcontentsline{toc}{chapter}{修订历史}

\begin{longtable}[l]{ | m{2cm} | m{1.5cm} | m{12cm} | }\hline

    \rowcolor{lightgray}
    日期 & 版本 & \hypertarget{Release notes}{发布说明}\\\hline
    \endfirsthead

    \multicolumn{3}{c}{\bfseries 接上页}\\\hline
    \rowcolor{lightgray}
    日期 & 版本 & 发布说明\\\hline
    \endhead

    \multicolumn{3}{r}{\textbf{见下页}}\\
    \endfoot

    \endlastfoot

    % % table contents
    % 2022-xx-xx     & v1.1  &
    % 更新以下章节：
    % \begin{itemize}
    %     \item 章节 \ref{mod:lowpowm} \textit{\nameref{mod:lowpowm}}
    %     \item 章节 \ref{mod:usbotg} \textit{\nameref{mod:usbotg}}
    % \end{itemize}
    % \\\hline
    2024- & v1.6 &
    更新以下章节：
        \begin{itemize}
            \item 章节 \ref{mod:lowpowm} \textit{\nameref{mod:lowpowm}}：
                \begin{itemize}
                \item 更新预设功耗模式相关的描述 %TRMC6-42193
                \item 增加表 \ref{tab:lpm-wakeup-source} 下方有关 EXT1 的说明 %DOC-8414
                \end{itemize}
            \item 章节 \ref{mod:iomuxgpio} \textit{\nameref{mod:iomuxgpio}}：更新 Hold 特性相关的描述 %DOC-7326
            \item 章节 \ref{mod:sensor} \textit{\nameref{mod:sensor}}：删除 SAR ADC2 测量内部电压/信号的相关描述 %DOC-7966
        \end{itemize}
    \\\hline
    2024-4-18 & v1.5 &
    更新以下章节：
        \begin{itemize}
            \item 章节 \ref{mod:dma} \textit{\nameref{mod:dma}}：更新 suc\_eof 和 EOF 标志的相关描述 %DOC-5478
            \item 章节 \ref{mod:iomuxgpio} \textit{\nameref{mod:iomuxgpio}}：更新 GPIO17 和 GPIO18 的驱动强度 %DOC-7383
            \item 章节 \ref{mod:digsig} \textit{\nameref{mod:digsig}}和章节 \ref{mod:bootctrl} \textit{\nameref{mod:bootctrl}}：修正笔误 %DOC-7472 %DOC-7629
            \item 章节 \ref{mod:uart} \textit{\nameref{mod:uart}}：更新清除 wake\_up 的说明 %DOC-7076
            \item 章节 \ref{mod:ledpwm} \textit{\nameref{mod:ledpwm}}：更新表 \ref{tab:ledc-common-freq-res} 中的最低精度  %DOC-7304
        \end{itemize}
    \\\hline
    2024-01-30     & v1.4  &
    更新以下章节：
    \begin{itemize}
        \item 章节 \ref{mod:pie} \textit{\nameref{mod:pie}}：在章节 \ref{chp:pie_instruction_simple_introduction} 增加数据交换指令列表；增加 LD.QR、ST.QR、以及 MV.QR 指令的描述；修正一处笔误 %DOC-6202；%DOC-5679
        \item 章节 \ref{mod:efuse} \textit{\nameref{mod:efuse}}：修改 \hyperref[fielddesc:EFUSEPINPOWERSELECTION]{EFUSE\_PIN\_POWER\_SELECTION} 寄存器的描述 % DOC-5591
        \item 章节 \ref{mod:iomuxgpio} \textit{\nameref{mod:iomuxgpio}}：从表格 \ref{tab:sysmem-base-address} 中删除辅助调试；更新章节 \ref{PadHold} 和表 \ref{tab:io_mux_reg_summary} 中的描述 %DOC-5892 %DOC-6733 %DOC-6563
        \item 章节 \ref{mod:bootctrl} \textit{\nameref{mod:bootctrl}}：在小节 \ref{sec:boot-mode-control} 中添加 SPI Download Boot 相关信息，并将 Download Boot 重命名为 Joint Download Boot；新增 EFUSE\_DIS\_FORCE\_DOWNLOAD 和 EFUSE\_DIS\_DOWNLOAD\_MODE 对芯片 Boot 控制的细节描述 %DOC-5923 %DOC-5912
        \item 章节 \ref{mod:intmtrx} \textit{\nameref{mod:intmtrx}}：删除 ASSIST\_DEBUG\_INTR 中断源 %DOC-6562
        \item 章节 \ref{mod:lowpowm} \textit{\nameref{mod:lowpowm}}：更新寄存器 \hyperref[fielddesc:RTCCNTLWDTWKEY]{RTC\_CNTL\_WDT\_WKEY} 的描述 %DOC-6760
        \item 章节 \ref{mod:timg} \textit{\nameref{mod:timg}}：更新 \hyperref[fielddesc:TIMGWDTCLKPRESCALE]{TIMG\_WDT\_CLK\_PRESCALE} 的描述 %DOC-5697
        \item 章节 \ref{mod:permctrl} \textit{\nameref{mod:permctrl}}：删除 Debug Assist 的访问权限配置  %DOC-6562
        \item 章节 \ref{mod:sysreg} \textit{\nameref{mod:sysreg}}：优化 \hyperref[fielddesc:SYSTEMCONTROLCORE1MESSAGE]{SYSTEM\_CONTROL\_CORE\_1\_MESSAGE} 寄存器域的描述 % DOC-5769
        \item 章节 \ref{mod:uart} \textit{\nameref{mod:uart}}：更新产生 wake\_up 所需的上升沿个数 %DOC-7076
        \item 章节 \ref{mod:i2c} \textit{\nameref{mod:i2c}}：更新 \hyperref[regdesc:I2CCOMD0REG]{I2C\_COMD0\_REG}、\hyperref[fielddesc:I2CSDAFORCEOUT]{I2C\_SDA\_FORCE\_OUT} 和 \hyperref[fielddesc:I2CSCLFORCEOUT]{I2C\_SCL\_FORCE\_OUT} 的描述 %DOC-6428, DOC-6511
        \item 章节 \ref{mod:mcpwm} \textit{\nameref{mod:mcpwm}}：增加递增递减模式下相关配置的说明 %DOC-7080
    \end{itemize}
    新增章节 \hyperref[interrupt-config-registers]{\textit{中断配置寄存器}}
    \\\hline
    2023-07-04 & v1.3 &
    更新以下章节：
    \begin{itemize}
        \item 章节 \ref{mod:ulp} \textit{\nameref{mod:ulp}}：增加 RTC I2C op\_code 的描述；在 \hyperref[regdesc:SENSSARI2CCTRLREG]{SENS\_SAR\_I2C\_CTRL\_REG} 寄存器中增加一些域；增加章节 “中断处理” %DOC-4933 DOC-5089
        \item 章节 \ref{mod:dma} \textit{\nameref{mod:dma}}：更新 GDMA\_IN\_SUC\_EOF\_CH\regindex{n}\_INT 中断和 GDMA\_INLINK\_DSCR\_ADDR\_CH\regindex{n} 字段的描述 % DOC-4730
        \item 章节 \ref{mod:sysmem} \textit{\nameref{mod:sysmem}}：更新表 \ref{tab:sysmem-base-address} 和图 \ref{fig:sysmem-address-mapping_structure} %DOC-5109
        \item 章节 \ref{mod:iomuxgpio} \textit{\nameref{mod:iomuxgpio}}：更新对 \hyperref[fielddesc:IOMUXPADPOWERCTRL]{IO\_MUX\_PAD\_POWER\_CTRL} 的描述 %DOC-5593
        \item 章节 \ref{mod:lowpowm} \textit{\nameref{mod:lowpowm}}：更新一处错误链接 %DOC-5116
        \item 章节 \ref{mod:permctrl} \textit{\nameref{mod:permctrl}} 和 \ref{mod:sysreg} \textit{\nameref{mod:sysreg}}：更新寄存器前缀 APB\_CTRL 为 SYSCON % DOC-5455
        \item 章节 \ref{mod:uart} \textit{\nameref{mod:uart}}：更新图 \nameref{fig:uart-arch-overview}、图 \nameref{fig:one-uart-structure}、图 \nameref{fig:uart-flow-control}；更新停止位最大位数和相关描述 %DOC-4826
        \item 章节 \ref{mod:usbotg} \textit{\nameref{mod:usbotg}}：增加对 USB OTG 寄存器获取方式的说明
        \item 章节 \ref{mod:usbserialjtag} \textit{\nameref{mod:usbserialjtag}}：在 \hyperref[fielddesc:USBSERIALJTAGPULLUPVALUE]{USB\_SERIAL\_JTAG\_PULLUP\_VALUE} 位的描述中增加具体的配置上拉电阻值
        \item 章节 \ref{mod:sensor} \textit{\nameref{mod:sensor}}：更新图 \ref{fig:pattern-tab-1} 和图 \ref{fig:pattern-tab-2} 的样式表编号
    \end{itemize}
    新增章节 \hyperref[programming-reserved-field]{\textit{如何配置寄存器的保留域}}
    \\\hline
    2023-03-02         & v1.2  &
    更新以下章节：
    \begin{itemize}
        \item 章节 \ref{mod:sysmem} \textit{\nameref{mod:sysmem}}：更新章节 \ref{other:sysmem-internal-ROM-1} \textit{\nameref{other:sysmem-internal-ROM-1}} 中 Internal ROM 1 的描述
        \item 章节 \ref{mod:lowpowm} \textit{\nameref{mod:lowpowm}}：删除 UART 作为拒绝睡眠原因 %DOC-4335 %DOC-3662
        \item 章节 \ref{mod:timg} \textit{\nameref{mod:timg}}：更新读取定时器值的步骤 %DOC-3634 %DOC-3480
        \item 章节 \ref{mod:sysreg} \textit{\nameref{mod:sysreg}}：更新寄存器 \newline \hyperref[regdesc:SYSTEMPERIPRSTEN0REG]{SYSTEM\_PERIP\_RST\_EN0\_REG} 和 \hyperref[regdesc:SYSTEMPERIPCLKEN0REG]{SYSTEM\_PERIP\_CLK\_EN0\_REG} 的描述 %DOC-4548
        \item 章节 \ref{mod:uart} \textit{\nameref{mod:uart}}：新增终止状态 (break condition) 的相关描述 %DOC-3633
        \item 章节 \ref{mod:usbserialjtag} \textit{\nameref{mod:usbserialjtag}}：新增 USB-OTG Download 下载模式下 IO 焊盘状态变化相关描述 %doc-4373
        \item 章节 \ref{mod:ledpwm} \textit{\nameref{mod:ledpwm}}：新增占空比精度计算公式和表 \textit{\nameref{tab:ledc-common-freq-res}}
        \item 章节 \ref{mod:sensor} \textit{\nameref{mod:sensor}}：\newline 新增 \hyperref[regdesc:APBSARADCAPBSARADC1DATASTATUSREG]{APB\_SARADC\_APB\_SARADC1\_DATA\_STATUS\_REG} 寄存器相关描述，更新 \hyperref[fielddesc:SENSFORCEXPDSAR]{SENS\_FORCE\_XPD\_SAR} 相关描述，并新增关于触摸传感器限制的说明
    \end{itemize}
    其他微小改动
    \\\hline
    2022-10-31         & v1.1  &
    更新以下章节：
    \begin{itemize}
        \item 章节 \ref{mod:pie} \textit{\nameref{mod:pie}}
        \item 章节 \ref{mod:efuse} \textit{\nameref{mod:efuse}}
        \item 章节 \ref{mod:sensor} \textit{\nameref{mod:sensor}}%DOC-3564
    \end{itemize}
    更新 \hyperref[glossary]{词汇列表} 小节
    \\\hline
    2022-09-08         & v1.0  &
    新增以下章节：
    \begin{itemize}
        \item 章节 \ref{mod:pie} \textit{\nameref{mod:pie}}
    \end{itemize}
    更新以下章节：
    \begin{itemize}
        \item 章节 \ref{mod:intmtrx} \textit{\nameref{mod:intmtrx}}
        \item 章节 \ref{mod:efuse} \textit{\nameref{mod:efuse}}
        \item 章节 \ref{mod:rng} \textit{\nameref{mod:rng}}
        \item 章节 \ref{mod:usbotg} \textit{\nameref{mod:usbotg}}
        \item 章节 \ref{mod:lowpowm} \textit{\nameref{mod:lowpowm}}
    \end{itemize}
    \\\hline
    2022-08-04     & v0.8  &
    新增以下章节：
    \begin{itemize}
        \item 章节 \ref{mod:i2s} \textit{\nameref{mod:i2s}}
        \item 章节 \ref{mod:spi} \textit{\nameref{mod:spi}}
    \end{itemize}
    更新以下章节：
    \begin{itemize}
        \item 章节 \ref{mod:efuse} \textit{\nameref{mod:efuse}} %doc-3121
        \item 章节 \ref{mod:lowpowm} \textit{\nameref{mod:lowpowm}} %register description about w1ts/w1tc
    \end{itemize}
    \\\hline
    2022-06-30     & v0.7  &
    新增以下章节：
    \begin{itemize}
        \item 章节 \ref{mod:sensor} \textit{\nameref{mod:sensor}}
        \item 章节 \ref{mod:lowpowm} \textit{\nameref{mod:lowpowm}}
    \end{itemize}
    更新以下章节：
    \begin{itemize}
        \item 章节 \ref{mod:sysmem} \textit{\nameref{mod:sysmem}}
        \item 更新时钟名称：
        \begin{itemize}
        \item FOSC\_CLK 更新为 RC\_FAST\_CLK
        \item FOSC\_DIV\_CLK 更新为 RC\_FAST\_DIV\_CLK
        \item RTC\_CLK 更新为 RC\_SLOW\_CLK
        \item SLOW\_CLK 更新为 RTC\_SLOW\_CLK
        \item FAST\_CLK 更新为 RTC\_FAST\_CLK
        \item PLL\_80M\_CLK 更新为 PLL\_F80M\_CLK
        \item PLL\_160M\_CLK 更新为 PLL\_F160M\_CLK
        \item PLL\_240M\_CLK 更新为 PLL\_D2\_CLK
    \end{itemize}%doc-3141
    \end{itemize}
    \\\hline
    2022-06-01     & v0.6  &
    更新以下章节：
    \begin{itemize}
        \item 章节 \ref{mod:dma} \textit{\nameref{mod:dma}}
        \item 章节 \ref{mod:efuse} \textit{\nameref{mod:efuse}} %DOC-2901
        \item 章节 \ref{mod:bootctrl} \textit{\nameref{mod:bootctrl}}
    \end{itemize}
    \\\hline
    2022-04-06     & v0.5  &
    新增以下章节：
    \begin{itemize}
        \item 章节 \ref{mod:permctrl} \textit{\nameref{mod:permctrl}}
        \item 章节 \ref{mod:wctrl} \textit{\nameref{mod:wctrl}}
    \end{itemize}
    更新以下章节：
    \begin{itemize}
        \item 章节 \ref{mod:sysmem} \textit{\nameref{mod:sysmem}}
    \end{itemize}
    \\\hline
    2022-02-24     & v0.4  &
    新增章节 \ref{mod:camlcdctrl} \textit{\nameref{mod:camlcdctrl}}
    更新以下章节：
    \begin{itemize}
        \item 章节 \ref{mod:ulp} \textit{\nameref{mod:ulp}} %DOC-2517, DOC-2516, DOC-2710, DOC-2611
        \item 章节 \ref{mod:efuse} \textit{\nameref{mod:efuse}}
        \item 章节 \ref{mod:iomuxgpio} \textit{\nameref{mod:iomuxgpio}} %DOC-2730
        \item 章节 \ref{mod:resclk} \textit{\nameref{mod:resclk}}
        \item 章节 \ref{mod:bootctrl} \textit{\nameref{mod:bootctrl}} %DOC-2569
        \item 章节 \ref{mod:usbotg} \textit{\nameref{mod:usbotg}}
    \end{itemize}
    \\\hline
    2021-12-16 & v0.3 &
    新增以下章节：
        \begin{itemize}
            \item 章节 \ref{mod:ulp} \textit{\nameref{mod:ulp}}
            \item 章节 \ref{mod:dma} \textit{\nameref{mod:dma}}
            \item 章节 \ref{mod:systimer} \textit{\nameref{mod:systimer}}
            \item 章节 \ref{mod:glitch} \textit{\nameref{mod:glitch}}
            \item 章节 \ref{mod:i2c} \textit{\nameref{mod:i2c}}
            \item 章节 \ref{mod:mcpwm} \textit{\nameref{mod:mcpwm}}
            \item 章节 \ref{mod:rmt} \textit{\nameref{mod:rmt}}
        \end{itemize}
    更新以下章节：
        \begin{itemize}
            \item 章节 \ref{mod:efuse} \textit{\nameref{mod:efuse}}
            \item 章节 \ref{mod:rsa} \textit{\nameref{mod:rsa}}
            \item 章节 \ref{mod:hmac} \textit{\nameref{mod:hmac}}
            \item 章节 \ref{mod:digsig} \textit{\nameref{mod:digsig}}
        \end{itemize}
    \\\hline
    2021-09-30      & v0.2 &
        新增以下章节：
            \begin{itemize}
                \item 章节 \ref{mod:sysreg} \textit{\nameref{mod:sysreg}}
                \item 章节 \ref{mod:hmac} \textit{\nameref{mod:hmac}}
                \item 章节 \ref{mod:extmemencr} \textit{\nameref{mod:extmemencr}}
                \item 章节 \ref{mod:uart} \textit{\nameref{mod:uart}}
                \item 章节 \ref{mod:usbserialjtag} \textit{\nameref{mod:usbserialjtag}}
            \end{itemize}

    更新以下章节：
            \begin{itemize}
                \item 章节 \ref{mod:efuse} \textit{\nameref{mod:efuse}}
                \item 章节 \ref{mod:twai} \textit{\nameref{mod:twai}}
            \end{itemize}
            \\\hline
    2021-07-09 & v0.1 & 首次预发布\\\hline

\end{longtable}


\clearpage
%%%%%%%%%%%%%%%%%%%%%%%%%%%%%%%%%%%%%%%%%%%%%%%%%%%%%%%%%%%%%%%%%%%%%%%%%%
%%% For complete instructions:                                         %%%
%%% 1. Go to https://www.overleaf.com/project/60e2ed1aefb70e595778b04b %%%
%%% 2. Compile readme-general.tex                                      %%%
%%% 3. See Section "Updating Disclaimer and Copyright Notice"          %%%
%%%%%%%%%%%%%%%%%%%%%%%%%%%%%%%%%%%%%%%%%%%%%%%%%%%%%%%%%%%%%%%%%%%%%%%%%%


\clearpage
\thispagestyle{empty}

\newcommand*{\myhref}[2]{%
\hypersetup{urlcolor=myred}\href{#1}{#2}}

\makeatletter
\setlength{\@fpbot}{0pt}
\makeatother

\newcolumntype{M}{@{}>{\columncolor{white}[0pt][0pt]}c}

\begin{table}
\def\arraystretch{1.2}
\begin{tabular}{c}
\setlength{\extrarowheight}{7pt}

\begin{minipage}[c]{16.7cm}
    \begin{tabular}{l}
        \hspace{-0.35cm}
        \includegraphics[width=6cm,valign=b,clip]{esp-logo-standard-horizontal.pdf}
        \vspace{0.6cm}
    \end{tabular}
\end{minipage} \\

\begin{minipage}[c]{16.7cm}
    \linespread{1.2}

    \textbf{\large 免责声明和版权公告}
    \vspace{0.4em}
    
    \small
    
    本文档中的信息，包括供参考的 URL 地址，如有变更，恕不另行通知。
    \vspace{0.4em}
    
    本文档可能引用了第三方的信息，所有引用的信息均为“按现状”提供，乐鑫不对信息的准确性、真实性做任何保证。
    \vspace{0.4em}
    
    乐鑫不对本文档的内容做任何保证，包括内容的适销性、是否适用于特定用途，也不提供任何其他乐鑫提案、规格书或样品在他处提到的任何保证。
    \vspace{0.4em}
    
    乐鑫不对本文档是否侵犯第三方权利做任何保证，也不对使用本文档内信息导致的任何侵犯知识产权的行为负责。本文档在此未以禁止反言或其他方式授予任何知识产权许可，不管是明示许可还是暗示许可。
    \vspace{0.4em}
    
    Wi-Fi 联盟成员标志归 Wi-Fi 联盟所有。蓝牙标志是 Bluetooth SIG 的注册商标。
    \vspace{0.4em}
    
    文档中提到的所有商标名称、商标和注册商标均属其各自所有者的财产，特此声明。
    \vspace{0.4em}
    
    \textbf{版权归 © \the\year{} 乐鑫信息科技（上海）股份有限公司。保留所有权利。}
    \vspace{0.4em}

    \textbf{\myhref{https://www.espressif.com/}{www.espressif.com}}
\end{minipage}

\end{tabular}
\end{table}


\end{document}
